\chapter{The Electron Is Named}

\section{Naming makes it measurable}

The residue splits into faces, and the faces are frame-dependent readings of one invariant. This chapter finds the first
face that is \emph{stable} across every frame --- a reading the corridor settles into and holds, the same in every viewing ---
and it finds that this stable reading is a named thing: the electron. The claim of the section is that the naming is not
decoration laid on afterward but the very act that makes the reading measurable. A face with no stable name is a viewpoint; a
face the counting forces into one box, and keeps there, is a particle. Naming is what turns the second into the first.

The mechanism is the pigeonhole, and in physics it wears the dress of representation theory. Ask what stable labels a
quantum system can carry, and the answer is Wigner's: the stable readings are the irreducible representations of the
symmetry group, and there are only so many of them at a given order. The state space, however endlessly the construction
climbs it, decomposes into a \emph{finite} list of irreducible types --- boxes labeled by the eigenvalues of the commuting
conserved quantities, the Casimirs. An unbounded tower of states sorted into finitely many boxes must, by the pigeonhole,
crowd many states into each box; and a box crowded with indistinguishable states, all carrying the same conserved labels,
\emph{is} a particle type. The electron is such a box: the label the counting forces a whole family of states to share.

The Dirac equation is where this box is exhibited exactly. Demand a first-order relativistic wave equation and the algebra
forces the wavefunction to carry four components and the operators to be the gamma matrices, $\gamma^\mu$, with
$\{\gamma^\mu, \gamma^\nu\} = 2 g^{\mu\nu}$~\cite{dirac1928}. The four components are not a modeling choice; they are forced
by the requirement that the equation be first-order and relativistic, exactly as the pigeonhole's boxes are forced by
finiteness. Out of that forcing fall the electron's stable labels: a definite spin, a definite charge, a definite mass-shell.
The equation does not \emph{posit} the electron and then describe it; it counts the components the algebra allows and finds
the electron sitting in one of the boxes the count forced open. Naming, here, is the reading-off of the forced box.

That the naming is measurable and not decorative is the whole point, and it is worth pressing. A quantity that varied
frame to frame --- like the split of a field into electric and magnetic parts --- cannot by itself name a particle, because
it is a viewpoint, not a thing. What names the electron are the labels every frame agrees on: the invariants, the Casimirs,
the charge and spin that no boost or relabeling alters. Those are measurable precisely because they are stable, and stable
precisely because they are the box the counting forced rather than the face a frame happened to show. The naming is the
selection of the invariant labels, and the invariant labels are the measurable ones. To name the electron is to have found
the reading that survives every viewing --- which is the same as to have found something the bench can measure at all.

In the two verbs, the naming is a funge the counting forces closed. The tower of states is funged into finitely many
boxes --- gathered by shared invariant labels --- and the pigeonhole guarantees the funge completes: with more states than
boxes, the gathering must pile many into one, and each crowded box is a named type. The electron is that closed funge, the
bag the counting could not help but fill. Its name is not a tange we impose from outside; it is the label the funge itself
carries once the counting has forced it shut. The next section reads which label, exactly, the electron's box carries --- and
finds it the stable, conserved reading the corridor holds in every frame.

\section{The stable reading}

The counting forced the electron into a box; this section reads the label on the box. The electron enters the construction
not as an arbitrary tag but as the \emph{stable} reading of the corridor --- the face that is conserved, that every frame
agrees on, that the loop returns unchanged. Its charge is a definite invariant, a single unit of the conserved flux, and the
sign of that unit is the electron's signature. The section's work is to show that ``stable'' and ``measurable'' and
``conserved'' are one property, and that the electron is the corridor's reading that has it.

Stability, in this register, is conservation, and conservation is Noether's theorem read backwards~\cite{noether1918}. A
reading is stable exactly when it is the conserved current of a symmetry: the charge $Q = \oint \star F$ is conserved because
the gauge symmetry is exact, $\partial_\mu J^\mu = 0$, and a conserved quantity is one the flow of the system does not alter.
So the electron's charge is stable not by fiat but because it is the Noether charge of the exact $U(1)$ symmetry --- the flux
through any closed surface around the electron is the same, in every frame, at every time. The stable reading is the
conserved one, and the conserved one is the measurable one, because a bench can only reliably read what does not drift while
it is being read.

The particular value of the electron's stable charge is a single unit of the conserved flux, carried with a definite sign,
and the construction reads that sign as the negative one. This is the corridor's stable face settling on the electron's
signature: one quantum of charge, and the minus that distinguishes the electron from its positive twin. The sign is not a
magnitude the instrument derives from first principles --- it is a label, a discrete two-valued reading, this box rather than
its mirror --- and the instrument treats it as exactly that: a stable, invariant, sign-valued reading, not a number fitted to
a scale. The electron is the box carrying one negative unit of the conserved flux, and that description is complete without
any real number in it. The charge is counted, not measured; it is $-1$ of a unit, an integer label on a forced box.

There is a discipline to observe here, and it is the same one the whole book keeps. The electron's stable labels --- its
integer charge, its half-integer spin --- are \emph{counted}, discrete, exact; and the book will not confuse them with the
one quantity that is genuinely a bracketed magnitude, the coupling. Charge is $-1$, a label the counting forces and the
instrument reads exactly. The coupling is an interval, floored, that the instrument reads only as a bracket. The two must
never be run together: naming the electron's charge is an exact count, and it lends no license to think the coupling's
magnitude is likewise exactly derivable. The chapter names an integer label with confidence precisely because integer labels
are the kind of thing counting delivers whole --- and it withholds that confidence, in advance, from the magnitude that is
not.

In the two verbs, the stable reading is the funge that the loop returns unchanged. The corridor's faces are many, and most
are tanges a frame selects and another frame unselects; the electron's charge is the one funge every frame shares --- the
covariant invariant the loop carries around and hands back identical. To name the electron is to funge the corridor down to
its conserved label and find that funge stable under every viewing. The unstable faces are viewpoints; the stable funge is a
particle. So the electron enters as the corridor's conserved reading, an integer label on a counting-forced box --- and the
last section shows why that box could only ever have been reached by counting, because the tower that holds it cannot be a
line and must close into a loop.

\section{Named by counting, not fiat}

The electron is a stable, conserved, integer-labeled reading. This closing section establishes the last thing about the
naming: that it is forced by \emph{counting} and could not have been installed by decree. The deep reason is a fact about the
shape of the state space --- the tower of the electron's states cannot be laid out as a line, and because it cannot, it
closes into a loop, and the closing is what forces the discrete label. The electron is named by counting because the
geometry of its states leaves no other way to name it.

The physics is the double cover, and it is one of the strangest true things in the subject. The rotations of space form the
group $SO(3)$, but the electron's states do not represent $SO(3)$; they represent its double cover $SU(2)$, in which a
rotation by $2\pi$ returns the electron not to itself but to \emph{minus} itself, and only a rotation by $4\pi$ brings it
home~\cite{pauli1927,dirac1928}. The state space is not a line you can lay flat; it is a space that has to be traversed twice
before it closes. That two-ness is not a choice --- it is forced by the algebra of the spin operators, $[S_i, S_j] =
i\hbar\,\epsilon_{ijk} S_k$, whose smallest nontrivial representation is two-dimensional and double-valued. The tower cannot
be a line, so it closes into a loop, and the loop takes $4\pi$ to close.

That the tower closes into a loop is exactly why the label is forced and cannot be decreed. A quantity laid out along an open
line could take any value; there is nothing to force it to be an integer, and a value assigned to it would be a fiat. But a
quantity defined on a \emph{loop} must be single-valued around the loop, and single-valuedness quantizes it: the phase must
return to itself after a full circuit, $\oint \mathrm{d}\phi \in 2\pi\mathbb{Z}$, which forces the winding to be an integer.
The electron's discrete labels --- its integer charge, its half-integer spin under the double cover --- are the winding
numbers of a loop that must close, and winding numbers are counted, never chosen. This is the difference between a name that
is measured and a name that is posited: the posited name sits on a line and could have been anything; the counted name is a
winding around a loop and could only have been what the closing forced.

The construction supplies the loop and its winding; the identification of the winding with the electron's specific quantum
numbers is the physical gauge's reading, marked as such. What the instrument derives is that a state space which cannot be a
line closes into a loop, and that a loop forces its labels to be counted windings. What the physics reads is that these
windings are the electron's charge and spin. The book keeps them apart, so that ``the electron is named by counting'' is
understood as a structural fact about loops and windings, not as a claim that the instrument derived the electron's
magnitudes. The labels are counted; the magnitudes --- when the book reaches one, at the coupling --- are bracketed; and the
counting that names the electron lends nothing to the bracketing that will floor the coupling.

In the two verbs, naming-by-counting is the funge a loop forces and a fiat could never close. On a line, the label would be a
free tange --- a value selected by decree, ungrounded. On a loop, the label is the funge the closing forces: the winding must
be an integer because the phase must return to itself, and the integer is gathered by the loop, not chosen by us. The
electron is named because its states funge onto a loop whose winding counts, and the count is the name. So the corridor's
stable face is the electron, a counting-forced winding on a loop that takes two turns to close --- and the next chapter takes
that loop and turns it, reading the turn itself as the operation that carries one face into another and charges angular
momentum for the passage.
